\section{Minimum Cost to Connect All Points}
\label{sec:graph-min-cost-connect}

\problemstatement{Given $n$ points on a 2D plane, the cost to connect
two points is their Manhattan distance
$\lvert x_1-x_2\rvert+\lvert y_1-y_2\rvert$. Find the minimum total cost
to connect all points so that there is exactly one path between any two.}

\begin{patternbox}
\emph{``Connect all nodes at minimum total cost, one path between
any two''} is the literal definition of a \textbf{Minimum Spanning
Tree}. The points form a complete graph (every pair is an edge weighted
by Manhattan distance); \textbf{Kruskal's} algorithm with Disjoint Set
selects the cheapest cycle-free edges until all points are joined.
\end{patternbox}

\subsection*{The complete graph and Kruskal's greedy choice}

Build every edge $(i,j)$ with weight equal to the Manhattan distance.
Kruskal's sorts all edges ascending and adds each one only if its two
endpoints are not already connected (checked in near-constant time with
Disjoint Set), which is precisely how it avoids cycles. Adding exactly
$n-1$ such edges spans all points at the lowest possible total cost.

\begin{ideabox}[Disjoint Set Makes ``Would This Edge Form a Cycle?'' O(1)]
Kruskal's correctness rests on always taking the globally cheapest
edge that connects two so-far-separate components. Disjoint Set answers
``are these endpoints already in the same component?'' in
$O(\alpha(n))$, turning the cycle check -- otherwise an expensive
traversal -- into an instant lookup.
\end{ideabox}

\subsection*{Algorithm}

\begin{enumerate}
  \item Form all $\binom{n}{2}$ edges $(\textit{weight}, i, j)$ with
        \textit{weight} the Manhattan distance; sort ascending.
  \item Create a Disjoint Set over the $n$ points.
  \item Scan edges in order; if the endpoints are in different sets,
        unite them and add the weight to the total. Stop once $n-1$
        edges have been added.
\end{enumerate}

\subsection*{C++ implementation}

\begin{lstlisting}
#include <bits/stdc++.h>
using namespace std;

class DisjointSet {
    vector<int> parent, size;
public:
    DisjointSet(int n) : parent(n), size(n, 1) {
        for (int i = 0; i < n; i++) parent[i] = i;
    }
    int find(int x) {
        if (parent[x] == x) return x;
        return parent[x] = find(parent[x]);
    }
    bool unite(int x, int y) {
        int rx = find(x), ry = find(y);
        if (rx == ry) return false;                // already connected
        if (size[rx] < size[ry]) swap(rx, ry);
        parent[ry] = rx; size[rx] += size[ry];
        return true;
    }
};

int minCostConnectPoints(vector<vector<int>>& points) {
    int n = points.size();
    vector<tuple<int,int,int>> edges;              // {weight, i, j}
    for (int i = 0; i < n; i++)
        for (int j = i + 1; j < n; j++) {
            int w = abs(points[i][0] - points[j][0]) +
                    abs(points[i][1] - points[j][1]);
            edges.push_back({w, i, j});
        }
    sort(edges.begin(), edges.end());

    DisjointSet ds(n);
    int total = 0, used = 0;
    for (auto& [w, i, j] : edges) {
        if (ds.unite(i, j)) {                      // adds only if no cycle
            total += w;
            if (++used == n - 1) break;
        }
    }
    return total;
}

int main() {
    vector<vector<int>> points = {{0,0},{2,2},{3,10},{5,2},{7,0}};
    cout << minCostConnectPoints(points) << endl;  // 20
    return 0;
}
\end{lstlisting}

\subsection*{Dry run}

\dryrun{Points $\{(0,0),(2,2),(3,10),(5,2),(7,0)\}$.}

\begin{itemize}
  \item Cheapest edges: $(0,0)\!-\!(2,2)=4$, $(2,2)\!-\!(5,2)=3$,
        $(5,2)\!-\!(7,0)=4$, $(5,2)\!-\!(3,10)=10$.
  \item Kruskal picks $3$ (connect $1,3$), $4$ (connect $0,1$), $4$
        (connect $3,4$), then $10$ (connect $2$) -- four edges for five
        points.
  \item Total $=3+4+4+10=\mathbf{20}$; every later edge would close a
        cycle and is skipped.
\end{itemize}

\subsection*{Complexity}

\begin{complexitybox}
\textbf{Time Complexity.} $O(n^2\log n)$ -- $O(n^2)$ edges dominate the
sort. \textbf{Space Complexity.} $O(n^2)$ for the edge list. (Prim's
with an adjacency-free $O(n^2)$ scan uses $O(n)$ space and can be
preferable for dense point sets.)
\end{complexitybox}

\begin{takeawaybox}
``Minimum total cost to connect everything'' is always an MST. Modeling
the points as a complete graph and running Kruskal's with Disjoint Set
reduces the problem to sort-then-union. When the graph is dense (as a
complete graph is), consider Prim's $O(n^2)$ variant to avoid materializing
every edge.
\end{takeawaybox}
